%=====================================================================
\chapter{Formula and Mathematics Reference}\label{app:formulas}
%=====================================================================
\normalfigures

Every formula used in this handout, with the chapter it comes from and a note on
when it applies. Suitable for use in an open-book examination.

\section{Notation}

Throughout: $n$ observations, $p$ features, $x_i$ the $i$-th observation, $y_i$ its
true value, $\hat{y}_i$ the prediction, $w$ the model parameters.

\section{Descriptive Statistics (Chapter 5)}

\begin{center}
\rowcolors{2}{white}{lightbg}
\begin{longtable}{@{}p{4.5cm}p{4.6cm}p{4.9cm}@{}}
\toprule
\rowcolor{primary!15}
\textbf{Quantity} & \textbf{Formula} & \textbf{Note} \\
\midrule
\endhead
Mean & $\bar{x}=\dfrac{1}{n}\sum_{i=1}^{n}x_i$ & Sensitive to outliers \\
Sample variance & $s^2=\dfrac{1}{n-1}\sum_{i}(x_i-\bar{x})^2$ & $n-1$ corrects the bias \\
Standard deviation & $s=\sqrt{s^2}$ & Same units as the data \\
Interquartile range & $\text{IQR}=Q_3-Q_1$ & Robust; use with the median \\
Outlier fences & $Q_1-1.5\,\text{IQR}$, $Q_3+1.5\,\text{IQR}$ & Flags \emph{candidates} only \\
$z$-score & $z=\dfrac{x-\bar{x}}{s}$ & Standardisation, Chapter 7 \\
Covariance & $\text{cov}(x,y)=\dfrac{1}{n-1}\sum_i(x_i-\bar{x})(y_i-\bar{y})$ & Units are the product \\
Pearson correlation & $r=\dfrac{\text{cov}(x,y)}{s_x s_y}$ & \textbf{Linear} association only \\
\bottomrule
\end{longtable}
\end{center}

\section{Probability (Chapter 6)}

\[
P(\bar{A}) = 1 - P(A)
\qquad
P(A\cup B)=P(A)+P(B)-P(A\cap B)
\]
\[
P(A\cap B)=P(A)\,P(B\mid A)
\qquad\text{and }=P(A)P(B)\ \text{\emph{only if independent}}
\]
\[
P(A \mid B) = \frac{P(A\cap B)}{P(B)}
\qquad\qquad
\boxed{\;P(A\mid B) = \frac{P(B\mid A)\,P(A)}{P(B)}\;}\quad\text{(Bayes)}
\]
\[
E[X]=\sum_k k\,P(X=k)
\qquad
\text{Var}(X)=E[(X-E[X])^2]
\]

\begin{conceptbox}[Central Limit Theorem]
For samples of size $n$ from any distribution with mean $\mu$ and standard
deviation $\sigma$, the sample mean is approximately
$\mathcal{N}\!\left(\mu,\ \sigma/\sqrt{n}\right)$. \textbf{Uncertainty falls as
$\sqrt{n}$: quadruple the data to halve the error.}
\end{conceptbox}

\section{Inference (Chapter 9)}

\[
\text{SE}(\bar{x}) = \frac{s}{\sqrt{n}}
\qquad
\text{95\% CI} = \bar{x} \pm 1.96\,\frac{s}{\sqrt{n}}
\]
\[
z = \frac{\bar{x}-\mu_0}{s/\sqrt{n}}
\qquad
\text{two proportions: } z = \frac{\hat{p}_1-\hat{p}_2}
{\sqrt{\hat{p}(1-\hat{p})\left(\frac{1}{n_1}+\frac{1}{n_2}\right)}}
\]
\[
\text{Cohen's } d = \frac{\bar{x}_1-\bar{x}_2}{s_{\text{pooled}}}
\qquad
\text{Bonferroni: } \alpha' = \frac{\alpha}{m}
\]
\[
\text{sample size for two proportions: } n \approx \frac{16\,\bar{p}(1-\bar{p})}{\delta^2}
\ \text{per group}
\]

\section{Linear Algebra and Optimisation (Chapter 4)}

\[
\|x\| = \sqrt{\textstyle\sum_j x_j^2}
\qquad
d(a,b) = \sqrt{\textstyle\sum_j (a_j-b_j)^2}
\qquad
a\cdot b = \textstyle\sum_j a_j b_j
\]
\[
\cos\theta = \frac{a\cdot b}{\|a\|\,\|b\|}
\qquad\qquad
\boxed{\;w \leftarrow w - \alpha\,\nabla L(w)\;}\quad\text{(gradient descent)}
\]

\section{Regression (Chapter 12)}

\[
\hat{y} = w_0 + \sum_{j=1}^{p} w_j x_j
\qquad
e_i = y_i - \hat{y}_i
\qquad
\text{least squares minimises } \textstyle\sum_i e_i^2
\]
\[
\text{MAE}=\frac{1}{n}\sum_i|e_i|
\qquad
\text{RMSE}=\sqrt{\frac{1}{n}\sum_i e_i^2}
\qquad
R^2 = 1-\frac{\sum_i e_i^2}{\sum_i (y_i-\bar{y})^2}
\]
\[
\text{Ridge: } \sum_i e_i^2 + \lambda\sum_j w_j^2
\qquad
\text{Lasso: } \sum_i e_i^2 + \lambda\sum_j |w_j|
\]

\section{Classification (Chapter 13)}

\[
\sigma(z)=\frac{1}{1+e^{-z}}
\qquad
\text{odds ratio for } w_j = e^{w_j}
\qquad
\text{Gini } G = 1-\sum_k p_k^2
\]
\[
\text{log loss} = -\frac{1}{n}\sum_i\left[y_i\log\hat{p}_i + (1-y_i)\log(1-\hat{p}_i)\right]
\]

\section{Evaluation (Chapter 14)}

\dsfigh{%
\begin{tikzpicture}[font=\small]
\node[rounded corners=3pt, fill=greenhl!16, draw=greenhl, minimum width=2.6cm,
      minimum height=1.1cm] at (0,0) {TP};
\node[rounded corners=3pt, fill=accent!13, draw=accent, minimum width=2.6cm,
      minimum height=1.1cm] at (2.8,0) {FP};
\node[rounded corners=3pt, fill=accent!13, draw=accent, minimum width=2.6cm,
      minimum height=1.1cm] at (0,-1.25) {FN};
\node[rounded corners=3pt, fill=greenhl!16, draw=greenhl, minimum width=2.6cm,
      minimum height=1.1cm] at (2.8,-1.25) {TN};
\node[font=\scriptsize, text=primary] at (0,0.85) {actually +};
\node[font=\scriptsize, text=primary] at (2.8,0.85) {actually $-$};
\node[font=\scriptsize, text=primary, rotate=90] at (-1.65,0) {said +};
\node[font=\scriptsize, text=primary, rotate=90] at (-1.65,-1.25) {said $-$};

\node[anchor=west, align=left] at (4.6,-0.6) {%
$\text{Accuracy}=\dfrac{\text{TP}+\text{TN}}{\text{TP}+\text{TN}+\text{FP}+\text{FN}}$\\[7pt]
$\text{Precision}=\dfrac{\text{TP}}{\text{TP}+\text{FP}}$
\qquad
$\text{Recall}=\dfrac{\text{TP}}{\text{TP}+\text{FN}}$\\[7pt]
$\text{Specificity}=\dfrac{\text{TN}}{\text{TN}+\text{FP}}$
\qquad
$F_1=2\,\dfrac{P\cdot R}{P+R}$};
\end{tikzpicture}}%
{The confusion matrix and every metric derived from it.}{formula-confusion}

\begin{alertbox}[The Two Rules Worth Memorising]
Under class imbalance, \textbf{accuracy is dominated by the majority class} ---
always compare against the majority-class baseline first. And \textbf{ROC-AUC
flatters rare-event detectors} because its $x$-axis divides by a very large number
of negatives; report PR-AUC instead.
\end{alertbox}

\section{Neural Networks (Chapter 17)}

\[
z = \sum_j w_j x_j + b
\qquad
a = f(z)
\qquad
\text{ReLU}(z)=\max(0,z)
\]
\[
\text{softmax}(z)_k = \frac{e^{z_k}}{\sum_j e^{z_j}}
\qquad
\text{parameters in a dense layer} = (\text{in}\times\text{out}) + \text{out}
\]

\section{Text (Chapter 20)}

\[
\text{tf-idf}(t,d) = \text{tf}(t,d)\times\log\frac{N}{\text{df}(t)}
\]

\section{Quick Decision Rules}

\rowcolors{2}{white}{defbg}
\begin{longtable}{@{}p{6.0cm}p{8.0cm}@{}}
\toprule
\rowcolor{primary!15}
\textbf{Situation} & \textbf{Rule} \\
\midrule
\endhead
Mean $\gg$ median & Right-skewed; report the median (Ch 5) \\
$|r| > 0.85$ between two features & Multicollinearity; keep one (Ch 8) \\
Positives $< 5\%$ of data & Use PR-AUC, not ROC-AUC (Ch 14) \\
Train score $\gg$ validation score & Overfitting; regularise or get more data (Ch 11) \\
Both scores low and equal & Underfitting; more capacity or better features (Ch 11) \\
Validation $>$ train score & Suspect a bug or a leak (Ch 11) \\
Accuracy above 99\% on a hard problem & Assume leakage until proven otherwise (Ch 7) \\
Data has a time index & Walk-forward validation, never random (Ch 19) \\
Many rows per entity & Split by entity, not by row (Ch 11) \\
Distance- or gradient-based method & Scale the features (Ch 4, 7) \\
Tree-based method & Scaling is unnecessary (Ch 15) \\
Base rates differ between groups & Calibration and demographic parity cannot both hold (Ch 28) \\
$n < 100$ & Prefer a simple model you can explain (Ch 11) \\
\bottomrule
\end{longtable}
